% Formalization support lemmas.
%
% NOT book material. These declarations were invented in this workspace to
% carry a Lean formalization and have no counterpart in the source text --
% they carry no \dcref. They live here, outside the book chapters, so the
% book blueprint stays timeless mathematics while their \leanok marks and
% Lean anchors are preserved.
%
% Do not add book mathematics to this file, and do not add formalization
% notes to the book chapters.

\begin{corollary}[Positive constant-curvature endpoints lie in the conjugate locus]
  \label{cor:dc-ch5-3-3-conjugate-locus-endpoint}
  \lean{Riemannian.Jacobi.mem_conjugateLocus_constCurvature_pos}
  \leanok
  \uses{cor:dc-ch5-3-3-first-conjugate,def:dc-ch5-3-4}
  Under the hypotheses of the preceding corollary, if $p=\gamma(0)$ then
  $\gamma(\pi/\sqrt{K_0})\in C(p)$. Thus the intrinsic formalized consequence
  of the sphere calculation is that every such first endpoint belongs to the
  conjugate locus; identifying it with the antipode, and proving
  $C(p)=\{-p\}$ for the round sphere, requires the additional global sphere
  geometry.
\end{corollary}

\begin{corollary}[First conjugate point in positive constant curvature]
  \label{cor:dc-ch5-3-3-first-conjugate}
  \lean{Riemannian.Jacobi.isFirstConjugatePointAt_constCurvature_pos}
  \leanok
  \uses{ex:dc-ch5-3-3,def:dc-ch5-3-1,lem:dc-ch8-2-1-no-conjugate}
  Let $M^n$ have constant sectional curvature $K_0>0$, with $n\geq 2$, and let
  $\gamma\colon[0,\pi/\sqrt{K_0}]\to M$ be a unit-speed geodesic. Then
  $\gamma(\pi/\sqrt{K_0})$ is the first conjugate point to $\gamma(0)$ along
  $\gamma$. Indeed, the sine Jacobi fields of \cref{ex:dc-ch5-3-3} give
  conjugacy at the endpoint, while the constant-curvature non-conjugacy
  criterion rules out every $0<t<\pi/\sqrt{K_0}$.
\end{corollary}

\begin{lemma}[pointwise Riemannian volume form]
  \label{lem:dc-ch1-2-11-pointwise}
  \uses{def:dc-ch1-2-1}
  \lean{Riemannian.DCPointwiseOrientation, Riemannian.DCVolumeFormAt,
    Riemannian.DCVolumeFormAt_eq_det}
  \leanok
  Let $(M,g)$ be an $n$-dimensional Riemannian manifold and choose an orientation
  $o_p$ of each tangent space.  At every $p\in M$ there is a canonical alternating
  $n$-form $\nu_p$ determined by $g_p$ and $o_p$.  If $(e_1,\ldots,e_n)$ is a
  positively oriented orthonormal basis of $T_pM$, then
  $$
  \nu_p(v_1,\ldots,v_n)=\det\bigl(\langle v_i,e_j\rangle_p\bigr)_{i,j}.
  $$
\end{lemma}
\begin{proof}
  \leanok
  The oriented inner-product space $(T_pM,g_p,o_p)$ has its canonical top
  alternating form.  By its defining invariance, this form is the determinant
  form of every positively oriented orthonormal basis, which gives the displayed
  formula.
\end{proof}

\begin{lemma}[synchronized finite chart partition for two curves]
  \label{lem:dc-ch8-2-1-chart-partition}
  \lean{Riemannian.Jacobi.exists_pair_curve_chart_partition}
  \leanok
  If $\gamma:[0,1]\to M$ and $\tilde\gamma:[0,1]\to\tilde M$ are continuous, then there are
  a monotone partition $0=\tau_0\leq\tau_1\leq\cdots$ which is eventually equal to $1$, and
  charts $\alpha_i$ on $M$ and $\tilde\alpha_i$ on $\tilde M$, such that each closed piece
  $[\tau_i,\tau_{i+1}]$ is mapped by $\gamma$ into the source of $\alpha_i$ and by
  $\tilde\gamma$ into the source of $\tilde\alpha_i$.  Thus both curves have a common finite
  subdivision by one-chart pieces.
\end{lemma}
\begin{proof}
  \leanok
  Apply the compactness lemma to the open cover of $[0,1]$ by the intersections of the pullbacks
  of chart sources along $\gamma$ and $\tilde\gamma$.
\end{proof}

\begin{lemma}[a common chart window for the Cartan transfer]
  \label{lem:dc-ch8-2-1-window-transfer}
  \lean{Riemannian.Jacobi.exists_common_chart_window_metricInner_mfderiv_eq_of_semiconjugacy_of_curvatureFormAt_of_isLocalDiffeomorphAt}
  \leanok
  \uses{lem:dc-ch8-2-1-window,
    lem:dc-ch8-2-1-exp-norm-transfer-general-discharged}
  In the setting of \cref{lem:dc-ch8-2-1-exp-norm-transfer-general-discharged}, suppose the two
  radial geodesics lie in fixed chart sources on $[0,1]$. Then there is one common interval
  $[a',b']$, with $a'<0<1<b'$, on which both chart-source conditions hold, and Cartan's
  curvature match for $\phi_t$ on that interval implies that the semiconjugacy preserves the
  metric at $q=\exp_p(v)$.
\end{lemma}
\begin{proof}
  \leanok
  \uses{lem:dc-ch8-2-1-window,
    lem:dc-ch8-2-1-exp-norm-transfer-general-discharged}
  Apply \cref{lem:dc-ch8-2-1-window} to each radial geodesic. If the resulting intervals are
  $[a_1,b_1]$ and $[a_2,b_2]$, take
  $$
  a'=\max\{a_1,a_2\},\qquad b'=\min\{b_1,b_2\}.
  $$
  Both strict endpoint inequalities persist, and the common interval lies in both original
  intervals. The two chart-source conditions therefore hold there, so the curvature match and
  \cref{lem:dc-ch8-2-1-exp-norm-transfer-general-discharged} give the metric identity. $\square$
\end{proof}

\begin{lemma}[cyclic polarization of the differentiated conformality equations]
  \label{lem:dc-ch8-5-differential-core}
  \uses{def:dc-ch8-5-conformal}
  \lean{Riemannian.liouville_secondDerivative_formula,
    Riemannian.liouville_secondDerivative_formula_of_inner_eq_zero}
  \leanok
  Let $E$ be a finite-dimensional real inner-product space. Let $L:E\to E$ be a
  conformal linear map with nonzero coefficient $\lambda$, let $B:E\times E\to E$
  be symmetric bilinear, and let $d\lambda:E\to\mathbb R$ have gradient
  $\nabla\lambda$. Suppose the differentiated conformality identity is
  \[
    \langle B(u,v),Lw\rangle+\langle Lv,B(u,w)\rangle
      =2\lambda\,d\lambda(u)\,\langle v,w\rangle
  \]
  for all $u,v,w$. Then
  \[
    B(u,v)=\frac1\lambda\bigl(d\lambda(u)Lv+d\lambda(v)Lu
      -\langle u,v\rangle L(\nabla\lambda)\bigr).
  \]
  In particular, if $u\perp v$, then
  \[
    B(u,v)=\frac1\lambda\bigl(d\lambda(u)Lv+d\lambda(v)Lu\bigr),
  \]
  which is equation (7) in the proof of Liouville's theorem.
\end{lemma}
\begin{proof}
  \leanok
  Write the differentiated identity for $(u,v,w)$, $(v,w,u)$, and $(w,u,v)$.
  Symmetry of $B$ and symmetry of the inner product identify the three repeated
  terms. Adding the first two equations and subtracting the third gives
  \[
    \langle B(u,v),Lw\rangle
      =\lambda\bigl(d\lambda(u)\langle v,w\rangle
        +d\lambda(v)\langle u,w\rangle
        -d\lambda(w)\langle u,v\rangle\bigr).
  \]
  The conformality identity for $L$ and
  $d\lambda(w)=\langle\nabla\lambda,w\rangle$ show that the right-hand side
  is the inner product of $Lw$ with the displayed formula for $B(u,v)$.
  Since $\lambda\ne0$, $L$ is injective and hence surjective in finite
  dimension; equality against every $Lw$ therefore gives equality of vectors.
  The orthogonal case is obtained by setting $\langle u,v\rangle=0$.
\end{proof}

\begin{lemma}[orientation determines the determinant of an orthogonal map]
  \label{lem:dc-ch9-3-7-orientation-det}
  \lean{Riemannian.Variation.det_sq_eq_one_of_isometry,
    Riemannian.Variation.det_eq_one_of_orientation_preserving,
    Riemannian.Variation.det_eq_neg_one_of_orientation_reversing}
  \leanok
  Let $V$ be a finite-dimensional oriented real inner-product space and let $A:V\to V$
  be an orthogonal transformation. Then $(\det A)^2=1$. If $A$ preserves the orientation,
  then $\det A=1$, while if $A$ reverses the orientation, then $\det A=-1$.
\end{lemma}
\begin{proof}
  \leanok
  In an orthonormal basis, the matrix of $A$ satisfies $A^{T}A=I$. Taking determinants gives
  $(\det A)^2=1$. Preservation of orientation is equivalent to $\det A>0$, and reversal of
  orientation is equivalent to $\det A<0$, which selects the corresponding root.
\end{proof}
